\documentclass[10pt,a4paper]{amsart}
\usepackage[margin=19mm]{geometry}
\usepackage{amsmath,amssymb,mathtools}
\usepackage[scr=rsfs,cal=euler]{mathalfa}
\usepackage{xfrac}
\usepackage[unicode,hidelinks,bookmarks=false]{hyperref}
\newcommand{\ie}{i.e.\ }
\newcommand{\cf}{cf.\ }
\newcommand{\resp}{respectively\ }
\newcommand{\Subsection}{\subsection}
\providecommand{\nobreakdash}{\nobreak\hbox{-}}
\setlength{\emergencystretch}{2em}
\multlinegap0pt
\usepackage{bm}
\usepackage{bbm}
\usepackage{dsfont}
\usepackage{xspace}
\usepackage{cancel}
\usepackage{mathtools}
\usepackage{amsthm}
\makeatletter
\@ifundefined{theorem}{\newtheorem{theorem}{Theorem}}{}
\makeatother
\def\BB{{\mathcal B}}
\def\PW{\text{\rm PW}}
\title[Etudes in the inverse spectral problem, II]{Etudes in the inverse spectral problem, II}
\author{Nikolai Makarov, Alexei Poltoratski}
\date{Source: arXiv:2509.10699v1 (CC BY 4.0)}
\begin{document}
\maketitle

\noindent\emph{Source and licence.} Etudes in the inverse spectral problem, II. Version arXiv:2509.10699v1 (CC BY 4.0), distributed under the Creative Commons Attribution 4.0 International licence, as recorded in the arXiv licence metadata for this exact version and shown on the abstract page https://arxiv.org/abs/2509.10699v1. Full rights notice: licenses/arxiv-2509.10699-CC-BY-4.0.txt.\par\smallskip

\noindent\textit{Editorial scope.} Counted unit: the Theorem whose source label is \texttt{thm2} in the article (source0.tex lines 709--709, source label \texttt{thm2}), delivered together with the article's own complete original proof of it (source0.tex lines 710--732, one \texttt{proof} environment of 23 lines). No mathematics is written, repaired or re-derived: statement and proof are the article's own text line for line. Required context retained uncounted: none. Every definition, display and label of those lines is retained unchanged. Omitted from this package and not redistributed: 1--708, 733--984. Nothing inside a retained excerpt is omitted or altered: every retained body is delivered exactly as the article's lines are written, so no omission receipt of any kind accompanies this package. Citations: the retained excerpts carry no citation command at all, so no bibliography entry of the article is transcribed and no reference is redistributed. Reference adaptation: none was needed and none was made; every reference inside the delivered unit resolves inside it, so no unresolved reference or citation is produced. Printed slips kept literal: no printed word, symbol, hypothesis or number of the counted statement or of its proof was altered. Layout and class: the article's own class is \texttt{amsart} with packages including amssymb, xcolor, hyperref; the wrapper is the lane's pinned \texttt{amsart} editorial wrapper at 10pt on A4, which restates the theorem-like environments the retained text needs and adds the article's own macro definitions that the retained text needs (\texttt{\textbackslash BB}, \texttt{\textbackslash PW}), with extra wrapper packages (bm, bbm, dsfont, xspace, cancel, mathtools). The delivered numbering is the unit's own. Assets: no retained span contains a figure environment, a graphics-inclusion command, a PDF-page inclusion, a table or a TikZ picture; no image file is copied into this package, the package declares no asset dependency and no image file is redistributed with it. Licence: version arXiv:2509.10699v1 of this article is distributed under the Creative Commons Attribution 4.0 International licence, as recorded in the arXiv licence metadata for this exact version and shown on the abstract page https://arxiv.org/abs/2509.10699v1. A full rights notice is delivered at licenses/arxiv-2509.10699-CC-BY-4.0.txt. Characters: every retained excerpt is pure ASCII, so the delivered engine, XeLaTeX with its default Latin Modern text font, reports no missing character.
\begin{theorem}\label{thm2}  $\mu\in(PW)$ is homogenous iff all its dB spaces are homogenous. \end{theorem}

\begin{proof} 
	
Suppose first that $\mu$ is homogeneous.	Let $F\in \BB_a$ and $0<t<1$. Then
$F\in \PW_a$ and 
$$G:=\sqrt t F(tz)\in\PW_{ta}\subset \PW_a=\BB_a$$
(last equation means "equal as sets"). It remains to show
$$\|G\|_{\BB_a}=
\|F\|_{\BB_a},$$  or equivalently
$$\int |G|^2~d\mu=\int |F|^2~d\mu.$$
We have
$$\int |G|^2~d\mu=\int tF(tx)^2~\rho(x)dx=\int F(y)^2~\rho(y)dy,$$
which means that all $\BB_a$ are homogeneous.

Suppose now that all $\BB_a$ are homogeneous. Let us show that $\mu(x)=\mu(tx)$ for all positive $t$.
If $F\in PW_a$ then
$$\int |F(x)|^2d\mu(x)=\int t|F(tx)|^2d\mu(x)=\int |F(y)|^2d\mu_t(y).$$
Hence, the measures $\mu(x)$ and $\mu_t(x)$ define the same norms on every $PW_a$ for any $c>0$.
Now, let us consider a det-normalized canonical system with the spectral measure $\mu$. 
By the definition of the spectral measure, $\mu_t(x)$ is also a spectral measure for the same system for any $t>0$.
From the uniqueness of the spectral measure, $\mu(x)=\mu_t(x)$.


\end{proof}

\end{document}
